\documentclass[a4paper,11pt,fontset=fandol]{ctexart}
\usepackage[margin=25mm]{geometry}
\usepackage{amsmath,graphicx}
\pagestyle{plain}
\input{assets/numbers.tex}
\begin{document}
\begin{center}
{\Large\bfseries 静态区间求和：前缀和实现与样例分析}
\end{center}
\section{问题与接口}
本作业针对不变数组上的多次区间求和。现有 Python 程序手写循环构建前缀数组。

\texttt{PrefixSum(values)} 接收数列。
\texttt{query(left, right)} 返回左闭右开区间
$[left,right)$ 中各元素的和。输入示例为 $[\InputValues]$。
本文说明递推关系、区间边界和重复查询时的取舍，并用给定的三个样例核对输出。
报告范围为已有实现与样例，不包含新的计时实验或程序修改。

\section{方法与边界}
\label{sec:method}
令数组长度为 $n$，前缀数组多保留一个初值，定义为
\begin{equation}
P[0]=0,\qquad P[k+1]=P[k]+a[k]\quad (0\leq k<n).
\label{eq:prefix}
\end{equation}
构造函数从 $[0]$ 出发，每读一个元素便将其与最后一个前缀值相加，再追加到列表。
因此 $P[k]$ 始终表示前 $k$ 个元素之和，列表最终有 $n+1$ 项。
查询时用
\begin{equation}
S(left,right)=P[right]-P[left]
\label{eq:query}
\end{equation}
消去 $left$ 之前的共同部分。式\eqref{eq:query}直接来自式\eqref{eq:prefix}的累计含义，
不是再次遍历查询区间。这个约定使 $right=n$ 对应列表中的最后一项，也使
$left=right$ 时两个相同前缀值相减，得到空区间和零。

实现先检查 $0\leq left\leq right\leq n$，不满足时抛出 \texttt{ValueError}。
这些边界条件来自现有代码的静态检查，本次样例执行没有覆盖异常分支。
对空输入，构造结果仍为 $[0]$；只允许区间 $[0,0)$。
这一行为可由代码推导，本次运行没有另设空输入测试。

从代码操作数看，预处理需要一次长度为 $n$ 的循环，时间和额外存储均为 $O(n)$；
构造之后，每次查询只做边界检查、两次取值和一次减法，在通常的整数运算模型下为 $O(1)$。
这是算法分析，不是实测耗时结论。相较每次直接遍历区间，该方案适合数组固定且查询较多的情形；
只有一次查询时，额外构造一份前缀数组未必值得。Python 大整数的运算成本还会受数值位数影响。

\clearpage
\section{样例结果与解释}
\label{sec:results}
实际执行现有程序后，得到前缀值 $[\PrefixValues]$。
图\ref{fig:prefix}画出相同输出中的累计值；横轴是已纳入的元素个数，纵轴是整数和，
用于观察式\eqref{eq:prefix}的递推过程，并非输入规模与耗时的关系。
\begin{figure}[htbp]
\centering
\includegraphics[width=.82\linewidth]{assets/prefix.pdf}
\caption{固定输入的前缀数组，各点直接来自已记录的程序输出。}
\label{fig:prefix}
\end{figure}

\begin{table}[htbp]
\centering
\input{assets/table.tex}
\caption{给定查询的独立答案与实际输出。}
\label{tab:queries}
\end{table}
表\ref{tab:queries}中的 $[0,4)$ 覆盖整个数组，输出为 $\FullSum$；
$[1,3)$ 只包含第二和第三个元素，使用 $P[3]-P[1]=12-2=\MiddleSum$；
$[2,2)$ 不包含元素，输出为 $\EmptySum$。三个查询分别检查右端点等于 $n$、中间子区间及空区间，
共 $\PassedCases/\TotalCases$ 个输出与任务材料中的独立答案一致。
数字宏、表格和图均从同一份归档程序输出派生，未手工修改某个图表中的数据。
这些结果支持所列样例的区间计算行为；负数、越界异常和较大输入尚未由本次运行验证，
因此不能把样例一致扩大为所有输入均正确，也不能据图评价运行速度。

\section{个人思考与限制}
\label{sec:thoughts}
我觉得统一成左闭右开的区间后，空区间可以直接用相同的两个前缀值相减，
不需要另加求和分支。前缀数组会多占一个元素和一份存储；如果只有一次查询，
直接遍历区间更简单。若数组经常修改，当前静态前缀和要重新构建，
是否换成其他结构还需要根据更新与查询的次数决定。

第\ref{sec:method}节的复杂度分析依赖静态数组与通常整数运算模型；
第\ref{sec:results}节的验证范围限于固定输入上的三个查询。
本报告没有测量耗时，也没有验证动态更新的数据结构。
\end{document}
